There are no additional level primes hidden in this argument.
Outside the prescribed support, $p$ is odd, $B_p$ is split, and
$R_p^0=\mathfrak{sl}_2(\mathbb Z_p)$ is self-dual for the bilinear
form of $Q_B=-a^2-bc$.  Its determinant is a unit.  Hence the
Fourier transform of $1_{R_p^0}$ equals itself; the unramified
Weil constant is one, as also follows from \cite[II.27]{weil} by integrating
over this self-dual lattice.  Upper unipotents, diagonal units,
and the Weyl element fix it.  These generate
$\operatorname{SL}_2(\mathbb Z_p)$, so the exponent there is zero.
At the finitely many prescribed primes the preceding $A_p$ suffice.
Conjugation by any $z\in\widehat B^\times$ is an isometry of $Q_B$,
commutes with the Weil representation, and preserves the
self-dual measure.  The same level and character therefore work
for every summand indexed by $z$.

\cite[III, Theorem 6]{weil} applies to this Schwartz--Bruhat function on
$V(\mathbb A)$ and proves invariance of its rational theta sum.
The normalized archimedean expression is
$\sum_x f(x)e^{2\pi i\tau Q_B(x)}$: the factor $y^{3/4}$ in the
Weil action is removed in passing to the classical form.
It has holomorphic weight $3/2$.  At a rational cusp, the slash
operation replaces the finite Schwartz function by its finite
Weil transform.  Its support meets $V(\mathbb Q)$ in a finite
union of rational lattice cosets, with bounded denominators.
Positive definiteness gives nonnegative rational exponents of
bounded denominator and normal convergence for $\Im\tau>0$.
This proves holomorphy at every cusp.  Theorem 6 alone asserts
automorphic invariance; this last argument supplies the
holomorphy conclusion used here.

